\begin{explanationbox}
\textbf{\textcolor{CExplanation}{一句话直觉}}
\\
三角形的内角和 $\pi$ 使正切函数通过和角公式产生对消，最终留下简洁的“和等于积”。

\medskip
\textbf{\textcolor{CExplanation}{核心拆解}}
\begin{description}[style=nextline, leftmargin=2.8em, labelsep=0.5em, itemsep=0.45em]
    \item[\textbf{要点一（内角代换）：}] 利用 $C = \pi - (A+B)$ 将 $\tan C$ 改写为 $-\tan(A+B)$。
    \item[\textbf{要点二（和角展开）：}] 展开 $\tan(A+B)$ 后通分，分子分母出现 $\tan A + \tan B$ 与 $\tan A \tan B$。
    \item[\textbf{要点三（乘积匹配）：}] 将 $\tan A + \tan B + \tan C$ 化为 $\dfrac{-\tan A \tan B (\tan A + \tan B)}{1 - \tan A \tan B}$，恰好等于 $\tan A \tan B \tan C$。
\end{description}

\medskip
\textbf{\textcolor{CExplanation}{几何本质（内角约束）}}
\\
三角形的内角和强制三个角中最多一个钝角，正切符号正好反映了这一事实：负的 $\tan$ 对应钝角，乘积符号直接与形状挂钩。

\medskip
\textbf{\textcolor{CExplanation}{代数意义（不变量关系）}}
\\
恒等式表明 $\tan A, \tan B, \tan C$ 不是独立的，它们满足一个对称有理关系，类似于三次方程韦达定理中的和与积相等的情形。

\medskip
\textbf{\textcolor{CExplanation}{考点价值}}
\begin{itemize}[leftmargin=*, itemsep=0.5em, topsep=0.4em]
    \item \textbf{考法一（知二求一）：} 已知两角正切，快速计算第三角正切，避免角的和差运算。
    \item \textbf{考法二（形状判定）：} 由三正切和的符号直接得出钝角三角形，减少计算量。
\end{itemize}

\medskip
\textbf{\textcolor{CExplanation}{顿悟点}}
\\
\textbf{把 $\tan C$ 替换为 $-\tan(A+B)$ 时，自然出现分母 $1-\tan A \tan B$，通分后 $\tan A \tan B$ 与和式完美重组，瞬间得到恒等式。}

\medskip
\textbf{\textcolor{CExplanation}{使用场景}}
\begin{itemize}[leftmargin=*, itemsep=0.5em, topsep=0.4em]
    \item \textbf{场景一（符号快速判形）：} 选择题中直接看三正切和的正负即知三角形有无钝角。
    \item \textbf{场景二（求值免去角度）：} 解析几何或三角求值题中，用乘积关系代替反三角运算。
\end{itemize}
\end{explanationbox}
